% Part: second-order-logic
% Chapter: syntax-and-semantics
% Section: expressive-power

\documentclass[../../../include/open-logic-section]{subfiles}

\begin{document}

\olfileid{sol}{syn}{exp}
\olsection{প্রকাশক্ষমতা}

\begin{explain}
দ্বিতীয়-ক্রমীয় চলরাশির উপর পরিমাণায়নের ফলে ভাষাটির প্রকাশক্ষমতা
প্রথম-ক্রমের যুক্তিবিদ্যার তুলনায় বিপুলভাবে বেড়ে যায়। দ্বিতীয়-ক্রমীয়
অস্তিত্বমূলক পরিমাণায়ন দিয়ে বলা যায়, বিশেষ কিছু ধর্মবিশিষ্ট অপেক্ষক বা
সম্পর্ক আছে। প্রথম-ক্রমের যুক্তিবিদ্যায় এটি বলার একমাত্র উপায় হলো ওই
কাজের জন্য একটি অ-যৌক্তিক সংকেত (অর্থাৎ, !!a{function} বা
!!{predicate}) নির্দিষ্ট করা। দ্বিতীয়-ক্রমীয় সার্বিক পরিমাণায়ন দিয়ে
বলা যায়, সংজ্ঞাক্ষেত্রের সব উপসেট, তার উপর সব সম্পর্ক, অথবা
!!{domain} থেকে !!{domain}-এ সব অপেক্ষকের কোনো ধর্ম আছে। প্রথম-ক্রমের
যুক্তিবিদ্যায় শুধু বলা যায়, ভাষার কোনো একটি অ-যৌক্তিক সংকেতে আরোপিত
উপসেট, সম্পর্ক বা অপেক্ষকের ওই ধর্ম আছে। আর যখন আমরা বলি কোনো
ধর্মবিশিষ্ট উপসেট, সম্পর্ক বা অপেক্ষক আছে, অথবা তাদের সবারই ওই ধর্ম আছে,
তখন ধর্মটি নির্দিষ্ট করতেও দ্বিতীয়-ক্রমীয় পরিমাণায়ন ব্যবহার করতে
পারি। ফলে প্রথম-ক্রমের যুক্তিবিদ্যায় সংজ্ঞায়িত করা যায় না এমন সম্পর্ক
সংজ্ঞায়িত করা এবং প্রথম-ক্রমের যুক্তিবিদ্যায় প্রকাশ করা যায় না এমন
সংজ্ঞাক্ষেত্রের ধর্ম প্রকাশ করা সম্ভব হয়।
\end{explain}

\begin{defn}
$\Struct{M}$ যদি ভাষা~$\Lang{L}$-এর জন্য !!a{structure} হয়, তবে
সম্পর্ক~$R \subseteq \Domain{M}^2$,~$\Lang{L}$-এ \emph{সংজ্ঞেয়}, যদি
শুধু $x$ ও~$y$ মুক্ত এমন কোনো !!{formula}~$!A_R(x, y)$ থাকে, যাতে
$R(a, b)$ সিদ্ধ হয় (অর্থাৎ, $\tuple{a,b} \in R$) যদি এবং কেবল যদি
$s(x) = a$ ও $s(y) = b$-এর জন্য $\Sat{M}{!A_R(x, y)}[s]$।
\end{defn}

\begin{ex}
প্রথম-ক্রমের যুক্তিবিদ্যায় সূত্র~$\eq[x][y]$ দিয়ে অভেদ
সম্পর্ক~$\Id{\Domain{M}}$ (অর্থাৎ, $\Setabs{\tuple{a,a}}{a \in
  \Domain{M}}$) সংজ্ঞায়িত করা যায়। দ্বিতীয়-ক্রমের যুক্তিবিদ্যায়
এই সম্পর্কটি \emph{$\eq$ ছাড়াই} সংজ্ঞায়িত করা যায়। কারণ $a$ ও~$b$
যদি~$\Domain{M}$-এর একই !!{element} হয়, তবে তারা~$\Domain{M}$-এর একই
উপসেটগুলির !!{element}s (যেহেতু সেট তার !!{element}s দ্বারা নির্ধারিত)।
বিপরীতভাবে, $a$ ও~$b$ ভিন্ন হলে তারা একই উপসেটগুলির !!{element}s নয়;
যেমন, $a \neq b$ হলে $a \in \{a\}$ কিন্তু $b \notin \{a\}$। তাই
“~$\Domain{M}$-এর একই উপসেটগুলির !!{element}s হওয়া” এমন একটি সম্পর্ক,
যা $a$ ও~$b$-এর ক্ষেত্রে সিদ্ধ হয় যদি এবং কেবল যদি $a = b$।
~$\Domain{M}$-এর সব উপসেটের উপর পরিমাণায়ন করা যায় বলে এই সম্পর্কটি
দ্বিতীয়-ক্রমের যুক্তিবিদ্যায় প্রকাশ করা যায়। অতএব নিচের !!{formula}
~$\Id{\Domain{M}}$-কে সংজ্ঞায়িত করে:
\[
\lforall[X][(X(x) \liff X(y))]
\]
\end{ex}

\begin{prob}
দেখাও যে $\lforall[X][(X(x) \lif X(y))]$ (লক্ষ করো:
$\lif$,~$\liff$ নয়!)~$\Id{\Domain{M}}$-কে সংজ্ঞায়িত করে।
\end{prob}

\begin{ex}
$R$ যদি একটি দ্বি-স্থানীয় !!{predicate} হয়, তবে $\Assign{R}{M}$ হলো
~$\Domain{M}$-এর উপর একটি দ্বি-স্থানীয় সম্পর্ক। কিছুটা বিভ্রান্তিকর
হতে পারে, তবু আমরা $R$-কে $R$-এর !!{predicate} এবং সম্পর্ক
~$\Assign{R}{M}$ উভয়ের জন্যই ব্যবহার করব। $R$-এর \emph{পরিযায়ী
আবরণ}~$R^*$ হলো সেই সম্পর্ক, যা $a$ ও~$b$-এর মধ্যে সিদ্ধ হয় যদি এবং
কেবল যদি কোনো $c_1$, \dots, $c_k$-এর জন্য $R(a,c_1)$, $R(c_1, c_2)$,
\dots, $R(c_k,b)$ সিদ্ধ হয়। এর মধ্যে $k = 0$-এর ক্ষেত্রটিও আছে;
অর্থাৎ, $R(a,b)$ সিদ্ধ হলে $R^*(a,b)$-ও সিদ্ধ হয়। তাই $R \subseteq
R^*$। বস্তুত, $R^*$ হলো $R$-কে অন্তর্ভুক্তকারী ক্ষুদ্রতম পরিযায়ী
সম্পর্ক। দ্বিতীয়-ক্রমের যুক্তিবিদ্যায় আমরা বলতে পারি, $X$ হলো
$R$-কে অন্তর্ভুক্তকারী একটি পরিযায়ী সম্পর্ক:
\begin{multline*}
  !B_R(X) \ident \lforall[x][\lforall[y][(R(x,y) \lif X(x, y))]] \land {}\\
\lforall[x][\lforall[y][\lforall[z][((X(x,y) \land X(y,z)) \lif X(x,
      z))]]].
\end{multline*}
প্রথম সংযোজ্যটি বলে $R \subseteq X$, আর দ্বিতীয়টি বলে $X$ পরিযায়ী।

$X$ এমন ক্ষুদ্রতম সম্পর্ক বলার অর্থ, $R$-কে অন্তর্ভুক্তকারী ও পরিযায়ী
প্রত্যেক সম্পর্কের মধ্যে সেটি নিজেই অন্তর্ভুক্ত। তাই নিচের !!{formula}
দিয়ে $R$-এর পরিযায়ী আবরণ সংজ্ঞায়িত করা যায়:
\[
R^*(X) \ident !B_R(X) \land \lforall[Y][(!B_R(Y) \lif
  \lforall[x][\lforall[y][(X(x, y) \lif Y(x,y))]])].
\]
$\Sat{M}{R^*(X)}[s]$ যদি এবং কেবল যদি $s(X) = R^*$। $R$-এর পরিযায়ী
আবরণ প্রথম-ক্রমের যুক্তিবিদ্যায় প্রকাশ করা যায় না।
\end{ex}

\end{document}
